\documentclass{report}
\usepackage{luatexja} % LuaTeXで日本語を使うためのパッケージ
\usepackage{luatexja-fontspec} % LuaTeX用の日本語フォント設定

% --- 数学関連 ---
\usepackage{amsmath, amssymb, amsfonts, mathtools, bm, amsthm} % 基本的な数学パッケージ
\usepackage{type1cm, upgreek} % 数式フォントとギリシャ文字
\usepackage{physics, mhchem} % 物理や化学の記号や式の表記を簡単にする

% --- 表関連 ---
\usepackage{multirow, longtable, tabularx, array, colortbl, dcolumn, diagbox} % 表のレイアウトを柔軟にする
\usepackage{tablefootnote, truthtable} % 表中に注釈を追加、真理値表
\usepackage{tabularray} % 高度な表組みレイアウト

% --- グラフィック関連 ---
\usepackage{tikz, graphicx} % 図の描画と画像の挿入
\usepackage{background} % ウォーターマークの設定
\usepackage{caption, subcaption} % 図や表のキャプション設定
\usepackage{float, here} % 図や表の位置指定

% --- レイアウトとページ設定 ---
\usepackage{fancyhdr} % ページヘッダー、フッター、余白の設定
\usepackage[top = 20truemm, bottom = 20truemm, left = 20truemm, right = 20truemm]{geometry}
\usepackage{fancybox, ascmac} % ボックスのデザイン

% --- 色とスタイル ---
\usepackage{xcolor, color, colortbl, tcolorbox} % 色とカラーボックス
\usepackage{listings, jvlisting} % コードの色付けとフォーマット

% --- 参考文献関連 ---
\usepackage{biblatex, usebib} % 参考文献の管理と挿入
\usepackage{url, hyperref} % URLとリンクの設定

% --- その他の便利なパッケージ ---
\usepackage{footmisc} % 脚注のカスタマイズ
\usepackage{multicol} % 複数段組
\usepackage{comment} % コメントアウトの拡張
\usepackage{siunitx} % 単位の表記
\usepackage{docmute}
% \usepackage{appendix}
% --- tcolorboxとtikzの設定 ---
\tcbuselibrary{theorems, breakable} % 定理のボックスと改ページ設定
\usetikzlibrary{decorations.markings, arrows.meta, calc} % tikzの装飾や矢印の設定

% --- 定理スタイルと数式設定 ---
\theoremstyle{definition} % 定義スタイル
\numberwithin{equation}{section} % 式番号をサブセクション単位でリセット

% --- hyperrefの設定 ---
\hypersetup{
  setpagesize = false,
  bookmarks = true,
  bookmarksdepth = tocdepth,
  bookmarksnumbered = true,
  colorlinks = false,
  pdftitle = {}, % PDFタイトル
  pdfsubject = {}, % PDFサブジェクト
  pdfauthor = {}, % PDF作者
  pdfkeywords = {} % PDFキーワード
}

% --- siunitxの設定 ---
\sisetup{
  table-format = 1.5, % 小数点以下の桁数
  table-number-alignment = center, % 数値の中央揃え
}

% --- 透かし画像の設定 --- 
\backgroundsetup{
  scale=0.5,                       % 画像のスケール
  % color=black,                   % 画像の色（透かし用に半透明が推奨）
  opacity=0.2,                   % 透かしの透明度（0が完全透明、1が完全不透明）
  angle=0,                       % 画像の角度
  position = current page.south east,  % ページの右下
  hshift=-6cm, % 右方向へのシフト（負の値で内側に移動）
  vshift=5cm,  % 上方向へのシフト（正の値で内側に移動）
  contents={\includegraphics{./fig/appilogo-circular-full.png}} % 画像のパス
}

% --- その他の設定 ---
\allowdisplaybreaks % 数式の途中改ページ許可
\newcolumntype{t}{!{\vrule width 0.1pt}} % 新しいカラムタイプ
\newcolumntype{b}{!{\vrule width 1.5pt}} % 太いカラム
\UseTblrLibrary{amsmath, booktabs, counter, diagbox, functional, hook, html, nameref, siunitx, varwidth, zref} % tabularrayのライブラリ
\setlength{\columnseprule}{0.4pt} % カラム区切り線の太さ
\captionsetup[figure]{font = bf} % 図のキャプションの太字設定
\captionsetup[table]{font = bf} % 表のキャプションの太字設定
\captionsetup[lstlisting]{font = bf} % コードのキャプションの太字設定
\captionsetup[subfigure]{font = bf, labelformat = simple} % サブ図のキャプション設定
\setcounter{secnumdepth}{4} % セクションの深さ設定
\newcolumntype{d}{D{.}{.}{5}} % 数値のカラム
\newcolumntype{M}[1]{>{\centering\arraybackslash}m{#1}} % センター揃えのカラム
\DeclareMathOperator{\diag}{diag}
\everymath{\displaystyle} % 数式のスタイル
\newcommand{\inner}[2]{\left\langle #1, #2 \right\rangle}
\renewcommand{\figurename}{図}
\renewcommand{\i}{\mathrm{i}} % 複素数単位i
\renewcommand{\laplacian}{\grad^2} % ラプラシアンの記号
\renewcommand{\thesubfigure}{(\alph{subfigure})} % サブ図の番号形式
\newcommand{\m}[3]{\multicolumn{#1}{#2}{#3}} % マルチカラムのショートカット
\renewcommand{\r}[1]{\mathrm{#1}} % mathrmのショートカット
\newcommand{\e}{\mathrm{e}} % 自然対数の底e
\newcommand{\Ef}{E_{\mathrm{F}}} % フェルミエネルギー
\renewcommand{\c}{\si{\degreeCelsius}} % 摂氏記号
\renewcommand{\d}{\r{d}} % d記号
\renewcommand{\t}[1]{\texttt{#1}} % タイプライタフォント
\newcommand{\kb}{k_{\mathrm{B}}} % ボルツマン定数
% \renewcommand{\phi}{\varphi} % ϕをφに変更
\renewcommand{\epsilon}{\varepsilon}
\newcommand{\fullref}[1]{\textbf{\ref{#1} \nameref{#1}}}
\newcommand{\reff}[1]{\textbf{図\ref{#1}}} % 図参照のショートカット
\newcommand{\reft}[1]{\textbf{表\ref{#1}}} % 表参照のショートカット
\newcommand{\refe}[1]{\textbf{式\eqref{#1}}} % 式参照のショートカット
\newcommand{\refp}[1]{\textbf{コード\ref{#1}}} % コード参照のショートカット
\renewcommand{\lstlistingname}{コード} % コードリストの名前
\renewcommand{\theequation}{\thesection.\arabic{equation}} % 式番号の形式
\renewcommand{\footrulewidth}{0.4pt} % フッターの線
\newcommand{\mar}[1]{\textcircled{\scriptsize #1}} % 丸囲み文字
\newcommand{\combination}[2]{{}_{#1} \mathrm{C}_{#2}} % 組み合わせ
\newcommand{\thline}{\noalign{\hrule height 0.1pt}} % 細い横線
\newcommand{\bhline}{\noalign{\hrule height 1.5pt}} % 太い横線

% --- カスタム色定義 ---
\definecolor{burgundy}{rgb}{0.5, 0.0, 0.13} % バーガンディ色
\definecolor{charcoal}{rgb}{0.21, 0.27, 0.31} % チャコール色
\definecolor{forest}{rgb}{0.0, 0.35, 0} % 森の緑色

% --- カスタム定理環境の定義 ---
\newtcbtheorem[number within = chapter]{myexc}{練習問題}{
  fonttitle = \gtfamily\sffamily\bfseries\upshape,
  colframe = forest,
  colback = forest!2!white,
  rightrule = 1pt,
  leftrule = 1pt,
  bottomrule = 2pt,
  colbacktitle = forest,
  theorem style = standard,
  breakable,
  arc = 0pt,
}{exc-ref}
\newtcbtheorem[number within = chapter]{myprop}{命題}{
  fonttitle = \gtfamily\sffamily\bfseries\upshape,
  colframe = blue!50!black,
  colback = blue!50!black!2!white,
  rightrule = 1pt,
  leftrule = 1pt,
  bottomrule = 2pt,
  colbacktitle = blue!50!black,
  theorem style = standard,
  breakable,
  arc = 0pt
}{proposition-ref}
\newtcbtheorem[number within = chapter]{myrem}{注意}{
  fonttitle = \gtfamily\sffamily\bfseries\upshape,
  colframe = yellow!20!black,
  colback = yellow!50,
  rightrule = 1pt,
  leftrule = 1pt,
  bottomrule = 2pt,
  colbacktitle = yellow!20!black,
  theorem style = standard,
  breakable,
  arc = 0pt
}{remark-ref}
\newtcbtheorem[number within = chapter]{myex}{例題}{
  fonttitle = \gtfamily\sffamily\bfseries\upshape,
  colframe = black,
  colback = white,
  rightrule = 1pt,
  leftrule = 1pt,
  bottomrule = 2pt,
  colbacktitle = black,
  theorem style = standard,
  breakable,
  arc = 0pt
}{example-ref}
\newtcbtheorem[number within = chapter]{exc}{Requirement}{myexc}{exc-ref}
\newcommand{\rqref}[1]{{\bfseries\sffamily 練習問題 \ref{exc-ref:#1}}}
\newtcbtheorem[number within = chapter]{definition}{Definition}{mydef}{definition-ref}
\newcommand{\dfref}[1]{{\bfseries\sffamily 定義 \ref{definition-ref:#1}}}
\newtcbtheorem[number within = chapter]{prop}{命題}{myprop}{proposition-ref}
\newcommand{\prref}[1]{{\bfseries\sffamily 命題 \ref{proposition-ref:#1}}}
\newtcbtheorem[number within = chapter]{rem}{注意}{myrem}{remark-ref}
\newcommand{\rmref}[1]{{\bfseries\sffamily 注意 \ref{remark-ref:#1}}}
\newtcbtheorem[number within = chapter]{ex}{例題}{myex}{example-ref}
\newcommand{\exref}[1]{{\bfseries\sffamily 例題 \ref{example-ref:#1}}}
% --- 再定義コマンド ---
% \mathtoolsset{showonlyrefs=true} % 必要な式番号のみ表示
\pagestyle{fancy} % ヘッダー・フッターのスタイル設定
\chead{応用量子物性講義ノート} % 中央ヘッダー
% \rhead{}
\fancyhead[R]{\rightmark}
\renewcommand{\sectionmark}[1]{\markright{\thesection\ #1}}
\cfoot{\thepage} % 中央フッターにページ番号
\lhead{}
\rfoot{Yuto Masuda and Haruki Aoki} % 右フッターに名前
\setcounter{tocdepth}{4} % 目次の深さ
\makeatletter
\@addtoreset{equation}{section} % サブセクションごとに式番号をリセット
\makeatother

% --- メタ情報 ---
\title{応用量子物性講義ノート}
\date{更新日\today}
\author{Yuto Masuda and Haruki Aoki}

\begin{document}
  特殊相対性理論(Special Relativity)は次の2つの事柄を原理とする．
  \begin{itembox}[l]{特殊相対性原理}
    あらゆる慣性系で同じ物理法則が成り立つ．
  \end{itembox}

  \begin{itembox}[l]{光速度不変の原理}
    あらゆる慣性形で真空中の光の速さは同一である．
  \end{itembox}

  この原理の下で成り立つ座標変換の法則(Lorentz変換)を導く．
  まず，慣性系$X$系の原点$O$とと$X'$系の原点$O'$が$t=t'=0$で一致している．$t=t'=0$で光が原点($O=O'$)を通過したとする．
  $X$系の空間座標を$(x,y,z)$，$X'$系の空間座標を$(x',y',z')$とすると，光速度不変の原理より，
  \begin{align}
    \frac{\sqrt{x^2 + y^2 + z^2}}{t} = \frac{\sqrt{x'^2 + y'^2 + z'^2}}{t'} = c
  \end{align}
  が成り立つ．上式から\textbf{世界長さ}(spacetime interval)
  \begin{align}
    s^2 = x^2 + y^2 + z^2 - (ct)^2
  \end{align}
  が不変量であることが導かれる．

  慣性系$X'$が$x$軸正の方向に速さ$v$で移動している．．このとき$y=y'，z=z'$である．わかりやすいように$T=it$とおく．光速不変より，
  \begin{align}
    x^2 - (ct)^2 &= x'^2 - (ct')^2\\
    x^2 + (cT)^2 &= x'^2 + (cT)^2
  \end{align}
  が成り立つ．これが回転座標変換と類似していることから，
  \begin{align}
    \begin{pmatrix}
      cT' \\ x'
    \end{pmatrix}
    =
    \begin{pmatrix}
      \cos\theta & -\sin\theta\\
      \sin\theta & \cos\theta
    \end{pmatrix}
    \begin{pmatrix}
      cT\\x
    \end{pmatrix}
  \end{align}
  と置く．表示を$t$に戻すと
  \begin{align}
    \begin{pmatrix}
      ct' \\ x'
    \end{pmatrix}
    =
    \begin{pmatrix}
      \cos\theta & \i\sin\theta\\
      \i\sin\theta & \cos\theta
    \end{pmatrix}
    \begin{pmatrix}
      ct\\x
    \end{pmatrix}
  \end{align}
  である．さらに，$\theta = \i\phi$とすると
  \begin{align}
    \begin{pmatrix}
      ct' \\ x'
    \end{pmatrix}
    =
    \begin{pmatrix}
      \cosh\phi & -\sinh\phi\\
      -\sinh\phi & \cosh\phi
    \end{pmatrix}
    \begin{pmatrix}
      ct\\x
    \end{pmatrix}
  \end{align}
  となる．よって，
  \begin{align}
    x' = (-\sinh\phi)ct + (\cosh\phi)x
  \end{align}
  を得る．$X$系において時刻$t$が経過したとする．$X$系から見るお$X'$系の原点の位置は$x=vt$である．一方，$X'$系から見ると$x'=0$である．
  よって上式から
  \begin{align}
    0 = (-\sinh\phi)ct + (\cosh\phi)vt
  \end{align}
  が成り立つ．よって，これを変形すると
  \begin{align}
    \frac{v}{c} = \frac{\sinh\phi}{\cosh\phi} = \tanh\phi
  \end{align}
  である．以上より，
  \begin{align}
    \begin{dcases}
      \sinh\phi = \frac{v/c}{\sqrt{1- (v/c)^2}}\\
      \cosh\phi = \frac{1}{\sqrt{1- (v/c)^2}}
    \end{dcases}
  \end{align}
  であることがわかる．したがって，
  \begin{itembox}[l]{Lorentz変換}
    \begin{align}
      \label{LorentzTransformation}
      \begin{pmatrix}
        ct' \\ x'
      \end{pmatrix}
      =
      \frac{1}{\sqrt{1 - (v/c)^2}}
      \begin{pmatrix}
        1 & -v/c\\
        -v/c & 1
      \end{pmatrix}
      \begin{pmatrix}
        ct\\x
      \end{pmatrix}
    \end{align}
  \end{itembox}
  を得る．Lorentz変換
  \begin{align}
    x' = \frac{x-vt}{\sqrt{1-(v/c)^2}}
  \end{align}
  において$v \ll c$とすると
  \begin{align}
    x' = x - vt
  \end{align}
  となる．これはGalilei変換と一致している．

  次に，相対論的効果を取り込んだ速度の合成則について示す．状況は，$X$系は静止し，$X'$系が速さ$v$で$x$軸方向に移動しているとする．さらに$X'$系では粒子が速さ$u'$で$x'$軸方向に運動している．
  $X$系から見た粒子の速度$V$を求める．式(\ref{LorentzTransformation})において$v\to -v$とすると
  \begin{align}
    \begin{pmatrix}
      c\dd{t} \\ \dd{x}
    \end{pmatrix}
    =
    \frac{1}{\sqrt{1 - (v/c)^2}}
    \begin{pmatrix}
      1 & v/c\\
      v/c & 1
    \end{pmatrix}
    \begin{pmatrix}
      c\dd{t'}\\\dd{x'}
    \end{pmatrix}
  \end{align}
  となる．$V=\dv{x}{t}$なので
  \begin{align}
    V = \dv{x}{t} &= \frac{v \dd{t'} + \dd{x'}}{\dd{t'} + (v/c^2)\dd{x'}}\\
    &= \frac{v + \dv{x'}{t'}}{1 + (v/c^2)\dv{x'}{t'}}\\
    &= \frac{v + u'}{1 + \frac{vu'}{c^2}}
  \end{align}
  を得る．
  \begin{itembox}[l]{速度の合成}
    \begin{align}
      V = \frac{v + u'}{1 + \frac{vu'}{c^2}}
    \end{align}
  \end{itembox}
  例として$u' = v = c$とすると
  \begin{align}
    V = \frac{2c}{1+c^2/c^2} = c
  \end{align}
  である．速度が合成されても光速を超えることは決してないことがわかる．

  \textbf{Lorentz収縮}について説明する．速さ$v$で運動している$X'$系から，$t'=0$において，静止している$X$系の2点を見る．1点は原点$O:(x,t)=(0,0)$
  もう1点は$P:(x,t)=(L,t)$とする．まず，原点$O$は$X'$系から見ると
  \begin{align}
    \begin{pmatrix}
      ct' \\ x'
    \end{pmatrix}
    =
    \frac{1}{\sqrt{1 - (v/c)^2}}
    \begin{pmatrix}
      1 & -v/c\\
      -v/c & 1
    \end{pmatrix}
    \begin{pmatrix}
      0\\0
    \end{pmatrix}
    =
    \begin{pmatrix}
      0\\0
    \end{pmatrix}
  \end{align}
  である．点$P$を$X'$系から見ると，
  \begin{align}
    \begin{pmatrix}
      ct' \\ x'
    \end{pmatrix}
    =
    \frac{1}{\sqrt{1 - (v/c)^2}}
    \begin{pmatrix}
      1 & -v/c\\
      -v/c & 1
    \end{pmatrix}
    \begin{pmatrix}
      ct\\L
    \end{pmatrix}
    =
    \frac{1}{\sqrt{1 - (v/c)^2}}
    \begin{pmatrix}
      ct - \frac{v}{c}L\\-vt + L
    \end{pmatrix}
  \end{align}
  である．$t'=0$で観測しているため，$t'=0$を代入し
  \begin{align}
    0 = ct - \frac{v}{c}L\\
    t = \frac{v}{c^2}L
  \end{align}
  を得る．よって，
  \begin{align}
    x' = \frac{1}{\sqrt{1 - (v/c)^2}} \qty(-\qty(\frac{v}{c})^2 L + L) = \sqrt{1-\qty(\frac{v}{c})^2}L
  \end{align}
  である．これは動いている慣性系から静止系での距離($L$)を測る($L'$)と縮んで見えることを意味している．
  \begin{itembox}[l]{Lorentz収縮}
    \begin{align}
      L' = \sqrt{1-\qty(\frac{v}{c})^2}L
    \end{align}
  \end{itembox}
  また，その対象に対して静止している観測者が測った距離を\textbf{固有長さ}(proper length)という．今回は$L$が固有長さである．

  次に\textbf{時間の遅れ}について説明する．同様の$X$系と$X'$系を考える．時計が$X'$系の原点$x'=0$に置かれており，観測者は$X'$系においてこの時計を見ている．
  $X'$系に置かれた時計の時刻が$t'$のときの時空点$P_1$，$t'+\Delta T_0$の時空点を$P_2$とする．$P_1$は，
  \begin{align}
    \begin{pmatrix}
      ct_1 \\ x_1
    \end{pmatrix}
    =
    \frac{1}{\sqrt{1 - (v/c)^2}}
    \begin{pmatrix}
      1 & v/c\\
      v/c & 1
    \end{pmatrix}
    \begin{pmatrix}
      ct'\\0
    \end{pmatrix}
    =
    \frac{1}{\sqrt{1 - (v/c)^2}}
    \begin{pmatrix}
      ct'\\vt'
    \end{pmatrix}
  \end{align}
  $P_2$は，
  \begin{align}
    \begin{pmatrix}
      ct_1 \\ x_1
    \end{pmatrix}
    =
    \frac{1}{\sqrt{1 - (v/c)^2}}
    \begin{pmatrix}
      1 & v/c\\
      v/c & 1
    \end{pmatrix}
    \begin{pmatrix}
      c(t' + \Delta T_0)\\0
    \end{pmatrix}
    =
    \frac{1}{\sqrt{1 - (v/c)^2}}
    \begin{pmatrix}
      c(t' + \Delta T_0)\\v(t' + \Delta T_0)
    \end{pmatrix}
  \end{align}
  である．よって，$X$系での時間経過$\Delta T = t_2 - t_1$は，
  \begin{align}
    \Delta T = t_2 - t_1 = \frac{1}{\sqrt{1 - (v/c)^2}} \Delta T_0
  \end{align}
  である．これは$X'$系は$X$系に比べて時間の流れが遅いことを示している．
  \begin{itembox}[l]{時間の遅れ}
    \begin{align}
      \Delta T = \frac{1}{\sqrt{1 - (v/c)^2}} \Delta \tau
    \end{align}
  \end{itembox}
  観測者に対して2つの事象が同一の空間座標で起きたとき，その時間間隔$\Delta \tau$を\textbf{固有時間}(proper time)という．

  電磁場の双対性について説明する．マクスウェル方程式は3次元では次の4つの式である．
  \begin{align}
    \begin{dcases}
      \nabla \cdot \bm{D} = \rho \\
      \nabla \cdot \bm{B} = 0\\
      \nabla \times \bm{E} = - \frac{\partial \bm{B}}{\partial t}\\
      \nabla \times \bm{H} = \frac{\partial \bm{D}}{\partial t} + \bm{j}
    \end{dcases}
  \end{align}
  これをLorentz変換に対して共変な形式に書き直す．まず，$\bm{B}$と$bm{E}$はベクトルポテンシャル$\bm{A}$とスカラーポテンシャル$\phi$を
  用いて
  \begin{align}
    \begin{dcases}
      \bm{E} = -\nabla \phi - \frac{\partial \bm{A}}{\partial t}\\
      \bm{B} = \nabla \times \bm{A}
    \end{dcases}
  \end{align}
  と表される．ここで，座標を
  \begin{align}
    (x^0, x^1, x^2, x^3) = (ct, x, y, z)
  \end{align}
  と表し，
  4元ベクトルポテンシャルを
  \begin{align}
    A^{\mu} &= (\phi, c\bm{A})\\
    A_{\mu} &= (\phi, -c\bm{A})
  \end{align}
  と定義する．この4元ベクトルポテンシャルを使うと電場と磁場は
  \begin{align}
    B_x &= c \frac{\partial A_z}{\partial y} - c \frac{\partial A_y}{\partial z}\\
    &= -\frac{\partial A_3}{\partial x^2} + \frac{\partial A_2}{\partial x^3}\\
    &= \partial_3 A_2 - \partial_2 A_3\\
    E_x &= -\frac{\partial \phi}{\partial x} - \frac{\partial A_x}{\partial t}\\
    &= -\frac{\partial A_0}{\partial x^1} + \frac{\partial A_1}{\partial x^0}\\
    &= \partial_0 A_1 - \partial_1 A_0
  \end{align}
  のように書ける．ここで，電磁場テンソル
  \begin{align}
    F_{\mu\nu} = \partial_{\mu}A_\nu - \partial_\nu A_\mu
  \end{align}
  を定義する．これを行列の形で表すと
  \begin{align}
    F_{\mu\nu}=
    \begin{pmatrix}
      F_{00} & F_{01} & F_{02} & F_{03}\\
      F_{10} & F_{11} & F_{12} & F_{13}\\
      F_{20} & F_{21} & F_{22} & F_{23}\\
      F_{30} & F_{31} & F_{32} & F_{33}
    \end{pmatrix}
    =
    \begin{pmatrix}
      0 & E_x & E_y & E_z \\
      -E_x & 0 & -cB_z & cB_y\\
      -E_y & cB_z & 0 & -cB_x\\
      -E_z& -cB_y & cB_x & 0
    \end{pmatrix}
  \end{align}
  である．明らかに$F_{\mu\nu} = F_{\nu\mu}$である．また，この定義から
  \begin{align}
    \partial_\mu F_{\nu\lambda} + \partial_\nu F_{\lambda\mu} + \partial_\lambda F_{\mu\nu} = 0
  \end{align}
  が成り立つ．実際に計算してみると
  \begin{align}
    &\partial_\mu F_{\nu\lambda} + \partial_\nu F_{\lambda\mu} + \partial_\lambda F_{\mu\nu} \\
    &= (\partial_\mu\partial_\nu A_\lambda - \partial_\mu\partial_\lambda A_\nu) + (\partial_\nu\partial_\lambda A_\mu - \partial_\nu\partial_\mu A_\lambda)
    + (\partial_\lambda\partial_\mu A_\nu - \partial_\lambda\partial_\nu A_\mu)\\
    &= 0
  \end{align}
  である．これはFaradayの電磁誘導の法則と磁束密度に関するGaussの法則を表している．
  例えば$\mu = 0,\nu = 1, \lambda =2$とすると
  \begin{align}
    &\partial_0 F_{12} + \partial_1 F_{20} + \partial_2 F_{01}\\
    &= -\frac{\partial (cB_z)}{\partial (ct)} + \frac{\partial}{\partial x}(-E_y) + \frac{\partial}{\partial y}(-E_x)\\
    &= -(\nabla\times \bm{E})_z - \frac{\partial B_z}{\partial t} = 0
  \end{align}
  $\mu=1,\nu=2,\lambda=3$とすると
  \begin{align}
    &\partial_1 F_{23} + \partial_2 F_{31} + \partial_3 F_{12}\\
    &= c(\partial_xB_x + \partial_yB_y + \partial_zB_Z)\\
    &= c\nabla\cdot\bm{B} = 0
  \end{align}
  が得られる．さらに，磁場/電束密度テンソル$H_{\mu\nu}$を
  \begin{align}
    H_{\mu\nu}=
    \begin{pmatrix}
      H_{00} & H_{01} & H_{02} & H_{03}\\
      H_{10} & H_{11} & H_{12} & H_{13}\\
      H_{20} & H_{21} & H_{22} & H_{23}\\
      H_{30} & H_{31} & H_{32} & H_{33}
    \end{pmatrix}
    =
    \begin{pmatrix}
      0 & cD_x & cD_y & cD_z \\
      -cD_x & 0 & -H_z & H_y\\
      -cD_y & H_z & 0 & -H_x\\
      -cD_z& -H_y & H_x & 0
    \end{pmatrix}
  \end{align}
  4次元の電流密度を$j^{\mu}$を
  \begin{align}
    j^{\mu} = (c\rho, \bm{j})
  \end{align}
  と定義する．明らかに$H_{\mu\nu} = -H_{\nu\mu}$である．
  すると，電束密度に関するGaussの法則とAmpèreの法則は次の式で表される．
  \begin{align}
    \partial^{\nu}H_{\nu\mu}=j_\mu
  \end{align}
  例えば$\mu=0$とすると
  \begin{align}
    \partial^1H_{10}+\partial^2H_{20} + \partial^3H_{30} &= j_0\\
    \partial_xD_x + \partial_yD_y + \partial_zD_z = \rho\\
    \nabla\cdot\bm{D} = \rho
  \end{align}
  が導かれる．さらに，電荷保存則
  \begin{align}
    \frac{\partial \rho}{\partial t} + \nabla\cdot\bm{j} = 0
  \end{align}
  から
  \begin{align}
    \frac{\partial (c\rho)}{\partial (ct)} + \nabla\cdot\bm{j} = 0
  \end{align}
  よって
  \begin{align}
    \partial_\mu j^\mu = 0
  \end{align}
  が得られる．以上をまとめるとMaxwell方程式は
  \begin{align}
    \begin{dcases}
      \partial_\mu F_{\nu\lambda} + \partial_\nu F_{\lambda\mu} + \partial_\lambda F_{\mu\nu} = 0\\
      \partial^{\nu}H_{\nu\mu}=j_\mu\\
      \partial_\mu j^\mu = 0
    \end{dcases}
  \end{align}
  と書くことができる．これらの式はテンソルとテンソル，ベクトルとベクトルというように，Lorentz変換に対して
  同じ変換則をもつものどうしが結ばれている．よってこれらはLorentz変換に対して共変である．
  \begin{itembox}[l]{Lorentz変換に対して共変なMaxwell方程式}
    \begin{align}
      &\partial_\mu F_{\nu\lambda} + \partial_\nu F_{\lambda\mu} + \partial_\lambda F_{\mu\nu} = 0\\
      &\partial^{\nu}H_{\nu\mu}=j_\mu\\
      &\partial_\mu j^\mu = 0
    \end{align}
  \end{itembox}

  運動する電荷と電流を例に電磁場の双対性を確認してみよう．導線には$z$軸上向きに電流$I$が流れている．電荷$+q$が$z$軸上向きに速さ$v_0$で運動している．
  導線の中では面電荷密度$\lambda_{\pm} = \pm\lambda/2$の正(負)電荷が速さ$v$で上(下)向きに動いているとする．
  まずは静止系で考える．導線の中では
  \begin{align}
    \lambda_{+} + \lambda_{-} = 0
  \end{align}
  が成り立つため，電気的に中性である．よってｍ導線の周りには電場は無い．電流は$I=\lambda v$と表される．
  導線の周りには
  \begin{align}
    \bm{B} = \frac{\mu_0I}{2\pi r} \bm{e}_{\phi} = \frac{\mu_0 \lambda v}{2\pi r} \bm{e}_{\phi}
  \end{align}
  の磁束密度が発生している．よって電荷は
  \begin{align}
    \bm{F} = q\bm{v}\times\bm{B} = -qv_0\frac{\mu_0 \lambda v}{2\pi r}\bm{e}_r
  \end{align}
  の力を感じる．

  これを電荷とともに動く系で考える．非相対論的に感がると$\bm{v}=0$であるため電荷は力を感じない．しかし，特殊相対性原理よりこれはありえない．
  相対論的効果を考慮してこの状況を眺める必要がある．
  観測者からは導線中の正電荷は$v-v_0$で，負電荷は$v+v_0$で運動して見える．それぞれでLorentz収縮を計算する．
  \begin{align}
    \lambda_{+} &= \frac{1}{\sqrt{1-\qty(\frac{v-v_0}{c})^2}}\frac{\lambda}{2}\\
    \lambda_{-} &= \frac{1}{\sqrt{1-\qty(\frac{v+v_0}{c})^2}}\qty(-\frac{\lambda}{2})
  \end{align}
  明らかに
  \begin{align}
    \lambda_{+} + \lambda_{-} \neq 0
  \end{align}
  であることがわかる．よって，導線の周りには電場が発生している．$v,v_0 \ll c$とすると
  \begin{align}
    \lambda_{+} + \lambda_{-} &= \frac{1}{\sqrt{1-\qty(\frac{v-v_0}{c})^2}}\frac{\lambda}{2} - \frac{1}{\sqrt{1-\qty(\frac{v+v_0}{c})^2}}\frac{\lambda}{2}\\
    &\simeq \qty(1 + \frac{1}{2}\qty(\frac{v-v_0}{c})^2)\frac{\lambda}{2} - \qty(1 + \frac{1}{2}\qty(\frac{v+v_0}{c})^2)\frac{\lambda}{2}\\
    &= - \frac{\lambda v_0 v}{c^2}
    &\equiv \Delta \lambda
  \end{align}
  と計算でき，導線は負に帯電していることがわかる．よって，電荷の受ける力は
  \begin{align}
    \bm{F} &= q\frac{\Delta \lambda}{2\pi r \epsilon_0}\bm{e}_r\\
    &=-qv_0\frac{\mu_0 \lambda v}{2\pi r}\bm{e}_r
  \end{align}
  これは先ほどの計算結果と一致している．以上の考察から，電場と磁場は観測する系によって入り混じることがわかる．
\end{document}